\documentclass[11pt]{article}
\usepackage[utf8]{inputenc}
\usepackage[T1]{fontenc}
\usepackage{amsmath,amssymb,amsthm,mathtools}
\usepackage[margin=1in]{geometry}
\usepackage{hyperref}

\theoremstyle{plain}
\newtheorem{theorem}{Theorem}
\newtheorem{corollary}[theorem]{Corollary}
\newtheorem{lemma}[theorem]{Lemma}

\theoremstyle{definition}
\newtheorem{definition}[theorem]{Definition}

\theoremstyle{remark}
\newtheorem*{remark}{Remark}

\DeclareMathOperator{\im}{Im}
\DeclareMathOperator{\re}{Re}

\title{The Probability Current of the Riemann Zeta Function\\
on the Critical Line}
\author{Stephen Crowley}
\date{April 2026}

\begin{document}
\maketitle

%==============================================================
\section{WKB Background}
%==============================================================

The WKB method (Wentzel, Kramers, Brillouin, 1926) is a technique
for constructing approximate solutions to the Schr\"odinger equation
in one spatial dimension when the potential varies slowly compared
to the de~Broglie wavelength. For the present purposes only the
structural content is needed, not the approximation.

\subsection*{The polar decomposition of a wavefunction}

A complex-valued function $\psi(x,t)$ satisfying the time-dependent
Schr\"odinger equation
\[
i\hbar\,\partial_t\psi = -\frac{\hbar^2}{2m}\,\partial_x^2\psi + V\psi
\]
can always be written in polar form
\[
\psi(x,t) = A(x,t)\,e^{iS(x,t)/\hbar},
\]
with $A \geq 0$ and $S$ real. Substituting and separating real and
imaginary parts yields two equations. The imaginary part gives the
\emph{continuity equation}
\[
\partial_t\rho + \partial_x j = 0,
\]
where the \emph{probability density} and \emph{probability current}
are
\[
\rho = A^2 = |\psi|^2,
\qquad
j = \frac{\hbar}{m}\,A^2\,\partial_x S
  = \frac{\hbar}{m}\,\im\!\left(\bar\psi\,\partial_x\psi\right).
\]
The real part gives a modified Hamilton--Jacobi equation for $S$.

\subsection*{What the current measures}

$\rho(x,t)\,dx$ is the probability of finding the particle in $[x,
x+dx]$ at time $t$. The current $j(x,t)$ is the net probability flux
crossing the point $x$ per unit time: positive means flux to the
right, negative to the left. The continuity equation states that
probability is locally conserved: any decrease in $\rho$ in a region
must equal the net outward flux through its boundary.

\subsection*{The velocity field}

The WKB velocity is $v = j/\rho = (\hbar/m)\,\partial_x S$.
In the classical limit $\hbar\to 0$ this coincides with the classical
particle velocity $p/m$ since $\partial_x S \to p$ (the action
gradient). The current is then $j = \rho v$, a density times a speed.

\subsection*{Units and normalisation}

In quantum mechanics $\int \rho\,dx = 1$ for a normalised state.
In the present paper $\psi$ is replaced by $\zeta(1/2+it)$, which is
not square-integrable on $[0,\infty)$, so $\rho$ is an \emph{unnormalised}
density. The role of $x$ is played by the height $t$, and the WKB
velocity field is $v(t) = \partial_t\arg\zeta(1/2+it)$. All
structural identities---the continuity equation, the relation $j =
\rho v$---hold without change.

%==============================================================
\section{Definitions}
%==============================================================

\begin{definition}[Standing objects]
\label{def:objects}
Let $\zeta(s)$ denote the Riemann zeta function. On the critical
line $\re(s) = 1/2$ define
\begin{align}
\mu(t) &\coloneqq \bigl|\zeta\bigl(\tfrac{1}{2}+it\bigr)\bigr|^{2}
&&\text{(local energy / probability density),}\\[4pt]
\vartheta(t) &\coloneqq
\im\log\Gamma\!\bigl(\tfrac{1}{4}+\tfrac{it}{2}\bigr) - \tfrac{t}{2}\log\pi
&&\text{(Riemann--Siegel theta),}\\[4pt]
Z(t) &\coloneqq e^{i\vartheta(t)}\,\zeta\!\bigl(\tfrac{1}{2}+it\bigr)
&&\text{(Hardy $Z$-function, real-valued for real $t$),}\\[4pt]
\Psi(t) &\coloneqq \int_0^t |\vartheta'(u)|\,\mu(u)\,du
&&\text{(cumulative phase clock),}\\[4pt]
\Psi'(t) &= |\vartheta'(t)|\,\mu(t)
&&\text{(cumulative phase density).}
\end{align}
The branch of $\log\Gamma$ is continuous on the upper half-plane and
real on the positive real axis. By Stirling,
\[
\vartheta'(t) = \tfrac{1}{2}\log\tfrac{t}{2\pi} + O(t^{-2}),
\]
so $|\vartheta'(t)| = \vartheta'(t)$ for all $t > t_c \approx 6.29$,
the unique zero of $\vartheta'$ on $(0,\infty)$. The absolute value
ensures $\Psi$ is strictly increasing on all of $[0,\infty)$.

The critical set is
\[
\mathcal{C} \coloneqq \bigl\{\,t\geq 0: \Psi'(t)=0\,\bigr\}
= \bigl\{\,t:\zeta\bigl(\tfrac{1}{2}+it\bigr)=0\,\bigr\} \cup \{t_c\}.
\]
$\mathcal{C}$ is discrete: zeros of $\zeta(\tfrac{1}{2}+it)$ are
isolated by the identity theorem for holomorphic functions, and
$\vartheta'$ has exactly one zero on $[0,\infty)$ by Stirling.
\end{definition}

\begin{definition}[Polar decomposition of $\zeta$]
\label{def:polar}
At every $t \notin \mathcal{C}$ write
\[
\zeta\bigl(\tfrac{1}{2}+it\bigr) = \sqrt{\mu(t)}\,e^{i\alpha(t)},
\qquad
\alpha(t) \coloneqq \arg\zeta\bigl(\tfrac{1}{2}+it\bigr),
\]
where $\arg$ denotes any fixed continuous branch of the argument on a
connected component of the complement of the zero set. At a zero $t_n$
the argument is undefined and the polar form breaks down; the primary
definitions below bypass this.
\end{definition}

\begin{definition}[Probability density and current]
\label{def:current}
The \emph{probability density} and \emph{probability current} of
$\zeta$ on the critical line are
\[
\rho(t) \coloneqq \mu(t)
= \bigl|\zeta\bigl(\tfrac{1}{2}+it\bigr)\bigr|^{2},
\]
\[
j(t) \coloneqq \im\!\Bigl(\,\overline{\zeta\bigl(\tfrac{1}{2}+it\bigr)}
\cdot\partial_t\,\zeta\bigl(\tfrac{1}{2}+it\bigr)\Bigr).
\]
These definitions require no branch choice and are valid for all
$t \in [0,\infty)$, including at every zero of $\zeta$.
\end{definition}

%==============================================================
\section{Main Results}
%==============================================================

\begin{theorem}[Explicit polar form of the current]
\label{thm:polar}
For all $t \notin \mathcal{C}$,
\[
j(t) = \mu(t)\,\alpha'(t) = \mu(t)\,\partial_t\arg\zeta\!\bigl(\tfrac{1}{2}+it\bigr).
\]
At every zero $t_n \in \mathcal{C}$ the definition gives $j(t_n)=0$
directly, consistently with $\lim_{t\to t_n} \mu(t)\alpha'(t) = 0$.
\end{theorem}

\begin{proof}
Write $\zeta_t = \sqrt{\mu(t)}\,e^{i\alpha(t)}$ for $t\notin\mathcal{C}$.
Differentiating,
\[
\partial_t\zeta_t
= \frac{\mu'(t)}{2\sqrt{\mu(t)}}\,e^{i\alpha(t)}
+ i\alpha'(t)\sqrt{\mu(t)}\,e^{i\alpha(t)}
= e^{i\alpha(t)}\!\left(\frac{\mu'(t)}{2\sqrt{\mu(t)}}
+ i\alpha'(t)\sqrt{\mu(t)}\right).
\]
Hence
\[
\bar\zeta_t\,\partial_t\zeta_t
= \sqrt{\mu(t)}\,e^{-i\alpha(t)}\cdot e^{i\alpha(t)}
\!\left(\frac{\mu'}{2\sqrt{\mu}}+i\alpha'\sqrt{\mu}\right)
= \frac{\mu'(t)}{2} + i\,\mu(t)\,\alpha'(t).
\]
Taking the imaginary part gives $j(t) = \mu(t)\,\alpha'(t)$.

At a zero $t_n$, Definition~\ref{def:current} gives
$j(t_n) = \im(0\cdot\partial_t\zeta_{t_n}) = 0$ since
$\zeta(t_n)=0$. Near a simple zero $\mu(t)\sim c(t-t_n)^2$ and
$|\alpha'(t)|$ is locally bounded by $|\vartheta'(t_n)| + O(1)$, so
$\mu(t)\,\alpha'(t)\to 0$. The two values are consistent.
\end{proof}

\begin{theorem}[Smooth part of the current equals $\Psi'$]
\label{thm:smooth}
Under the Riemann--Siegel factorisation $\zeta(\tfrac{1}{2}+it)
= e^{-i\vartheta(t)}Z(t)$ with $Z(t)$ real-valued, the phase
decomposes as
\[
\alpha(t) = -\vartheta(t) + \arg Z(t),\qquad t\notin\mathcal{C}.
\]
The full distributional derivative is
\[
\alpha'(t)
= -\vartheta'(t) + \partial_t\arg Z(t),
\]
where $\partial_t\arg Z$ contributes $\pm\pi\,\delta(t-t_n)$ at each
real zero $t_n$ of $Z$ (i.e.\ each critical-line zero of $\zeta$).
The current is therefore
\[
j(t) = -\mu(t)\,\vartheta'(t) + \mu(t)\,\partial_t\arg Z(t).
\]
The smooth leading term is
\[
j_{\mathrm{smooth}}(t) = -\mu(t)\,\vartheta'(t),
\qquad
|j_{\mathrm{smooth}}(t)| = \mu(t)|\vartheta'(t)| = \Psi'(t).
\]
\end{theorem}

\begin{proof}
From Definition~\ref{def:objects}, $Z(t) = e^{i\vartheta(t)}\zeta(\tfrac{1}{2}+it)$,
so $\arg\zeta = -\vartheta + \arg Z$. Differentiating: $\alpha' =
-\vartheta' + \partial_t\arg Z$.

Since $Z(t)$ is real-valued and $C^1$ between zeros, $\arg Z(t) \in
\{0,\pi\}$ between consecutive zeros and jumps by $\pm\pi$ at each
zero $t_n$. In the distributional sense,
$\partial_t\arg Z(t) = \sum_n (\pm\pi)\,\delta(t-t_n)$, with
sign determined by the direction of sign change of $Z$ at $t_n$.

Substituting into $j = \mu\alpha'$:
\[
j(t) = -\mu(t)\,\vartheta'(t) + \mu(t)\sum_n(\pm\pi)\,\delta(t-t_n).
\]
For each $n$, the distributional product $\mu(t)\,\delta(t-t_n)$ is
evaluated by pairing with a test function $\phi$:
\[
\int \mu(t)\,\delta(t-t_n)\,\phi(t)\,dt = \mu(t_n)\,\phi(t_n) = 0,
\]
since $\mu(t_n) = |\zeta(\tfrac{1}{2}+it_n)|^2 = 0$. Therefore the
distributional terms vanish identically and
$j(t) = -\mu(t)\,\vartheta'(t) = j_{\mathrm{smooth}}(t)$ holds
everywhere as a distribution. Taking the absolute value and using
Definition~\ref{def:objects},
\[
|j_{\mathrm{smooth}}(t)| = \mu(t)|\vartheta'(t)| = \Psi'(t).\qedhere
\]
\end{proof}

\begin{remark}
The conclusion $j = j_{\mathrm{smooth}}$ everywhere is exact, not
approximate. The branch-jump terms in $\partial_t\arg Z$ are real
features of $\alpha'$ but are annihilated by $\mu$ vanishing at each
zero. No approximation is involved.
\end{remark}

\begin{theorem}[Continuity equation in $\Psi$-time]
\label{thm:continuity}
Define $T = \Psi(t)$. In $T$-coordinates the density and current are
\[
\rho_T(T) \coloneqq \frac{\rho(\Psi^{-1}(T))}{\Psi'(\Psi^{-1}(T))}
= \frac{1}{|\vartheta'(\Psi^{-1}(T))|},
\qquad
j_T(T) \coloneqq \frac{j(\Psi^{-1}(T))}{\Psi'(\Psi^{-1}(T))}
= -\operatorname{sgn}(\vartheta'(\Psi^{-1}(T))).
\]
On the tail $t > t_c$, where $\vartheta'(t) > 0$, one has $j_T = -1$
(constant). The continuity equation
\[
\partial_T\rho_T + \partial_T j_T = 0
\]
holds identically.
\end{theorem}

\begin{proof}
For any Borel set $A \subset [0,\infty)$ the change of variables
$T = \Psi(t)$, $dT = \Psi'(t)\,dt$ gives
\[
\int_{\Psi(A)}\rho_T(T)\,dT = \int_A \rho(t)\,dt,
\]
so $\rho_T(T) = \rho(t)/\Psi'(t)$ at $T = \Psi(t)$. Substituting
$\rho = \mu$ and $\Psi' = \mu|\vartheta'|$:
\[
\rho_T = \frac{\mu}{\mu|\vartheta'|} = \frac{1}{|\vartheta'|}.
\]
Similarly $j_T = j_{\mathrm{smooth}}/\Psi' = -\mu\vartheta'/(\mu|\vartheta'|) =
-\operatorname{sgn}(\vartheta')$. On the tail, $\operatorname{sgn}(\vartheta')=1$,
so $j_T = -1$. Since $-1$ is constant, $\partial_T j_T = 0$, and
$\partial_T\rho_T = -\partial_T j_T = 0$ as well. The continuity
equation holds identically.
\end{proof}

\begin{corollary}[Unit-speed flow]
The time change $T = \Psi(t)$ is the unique reparametrisation of $t$
under which the WKB velocity $v_T = j_T/\rho_T$ satisfies
\[
v_T = \frac{j_T}{\rho_T}
= \frac{-1}{1/|\vartheta'|} = -|\vartheta'(\Psi^{-1}(T))|,
\]
which equals $-1$ in the sense that in $T$-time the current is
constant (unit magnitude). The non-constant speed $|\vartheta'(t)|
\sim \tfrac{1}{2}\log(t/2\pi)$ of the original $t$-parametrisation
is absorbed entirely into the clock $\Psi$.
\end{corollary}

\begin{theorem}[Total flux, $\Psi$-clock, and zero count]
\label{thm:flux}
The total flux of $|j_{\mathrm{smooth}}|$ over $[0,T]$ is
\[
\int_0^T |j_{\mathrm{smooth}}(t)|\,dt
= \int_0^T \Psi'(t)\,dt = \Psi(T).
\]
The smooth zero-counting function $\bar N(t) = \vartheta(t)/\pi + 1$
satisfies $\bar N'(t) = \vartheta'(t)/\pi$, so on the tail,
\[
|j_{\mathrm{smooth}}(t)| = \pi\,\mu(t)\,\bar N'(t),
\]
and therefore
\[
\Psi(T) = \pi\int_0^T \mu(t)\,\bar N'(t)\,dt.
\]
The total current flux, the elapsed $\Psi$-clock, and the
density-weighted zero count are the same integral.
\end{theorem}

\begin{proof}
The first equality is the definition of $\Psi(T) =
\int_0^T\Psi'(t)\,dt$ combined with $|j_{\mathrm{smooth}}| = \Psi'$
from Theorem~\ref{thm:smooth}. The second uses $\bar N'(t) =
\vartheta'(t)/\pi$ and $|j_{\mathrm{smooth}}| = \mu|\vartheta'| =
\pi\mu\bar N'$.
\end{proof}

\begin{theorem}[Radon--Nikodym derivative as current ratio]
\label{thm:RN}
Let $\gamma$ be the Borel measure $d\gamma(t) = \Psi'(t)\,dt =
|j_{\mathrm{smooth}}(t)|\,dt$. For $h\in\mathbb{R}$ let
$(T_h\gamma)(A) = \gamma(A+h)$. Then
\[
\frac{d(T_h\gamma)}{d\gamma}(t)
= \frac{\Psi'(t-h)}{\Psi'(t)}
= \frac{|j_{\mathrm{smooth}}(t-h)|}{|j_{\mathrm{smooth}}(t)|},
\]
for $\gamma$-almost every $t$. The measures $\gamma$ and $T_h\gamma$
are equivalent (mutually absolutely continuous).
\end{theorem}

\begin{proof}
For any Borel set $A$,
\[
(T_h\gamma)(A) = \gamma(A+h) = \int_{A+h}\Psi'(t)\,dt
= \int_A \Psi'(s+h)\,ds.
\]
So $T_h\gamma$ has Lebesgue density $\Psi'(t-h)$. Both $\Psi'(t)$
and $\Psi'(t-h)$ vanish only on $\mathcal{C}$ and $\mathcal{C}+h$
respectively. As established in Definition~\ref{def:objects},
$\mathcal{C}$ is discrete, hence countable, hence Lebesgue-null,
hence $\gamma$-null. The same holds for $\mathcal{C}+h$. Therefore
$\gamma(\{\Psi'=0\}) = 0$ and $T_h\gamma(\{\Psi'(\cdot-h)=0\}) = 0$,
so the two measures share all null sets. The chain rule
$d(T_h\gamma)/d\gamma = \Psi'(t-h)/\Psi'(t)$ follows, and is
positive $\gamma$-a.e. The final equality restates $\Psi' =
|j_{\mathrm{smooth}}|$ from Theorem~\ref{thm:smooth}.
\end{proof}

\begin{remark}
The Radon--Nikodym derivative is the ratio of current densities at
$t$ and $t-h$. When $h$ lies in the Cameron--Martin space $H_\gamma$
(the space of shifts for which this ratio is square-integrable with
respect to $\gamma$), the shift is absolutely continuous with respect
to $\gamma$. When $h\notin H_\gamma$ the translated measure
concentrates on a set disjoint from the support of $\gamma$ in the
measure-theoretic sense, and no Radon--Nikodym derivative exists.
\end{remark}

%==============================================================
\section{The Priestley Representation}
%==============================================================

The stationary pullback $X(T) = \zeta(\tfrac{1}{2}+i\Psi^{-1}(T))
/\sqrt{\Psi'(\Psi^{-1}(T))}$ has Fourier support in $[-2,0]$ and
admits the spectral integral
\[
X(T) = \int_{[-2,0]} e^{i\omega T}\,d\Phi(\omega),
\]
where $\Phi$ is a deterministic, convergent Borel measure on $[-2,0]$
(not a random measure, not a distribution). The recovery formula
$\zeta(\tfrac{1}{2}+it) = \sqrt{\Psi'(t)}\cdot X(\Psi(t))$ then
gives the Priestley oscillatory representation
\[
\zeta\!\bigl(\tfrac{1}{2}+it\bigr)
= \int_{[-2,0]}\sqrt{\Psi'(t)}\;e^{i\omega\Psi(t)}\,d\Phi(\omega).
\]
Expressed through $\Phi$, the density and current are
\begin{align}
\rho(t) &= \Psi'(t)\,\Bigl|\int_{[-2,0]}e^{i\omega\Psi(t)}\,d\Phi(\omega)\Bigr|^2,\\[6pt]
j(t) &= -\mu(t)\,\vartheta'(t)
= -\Psi'(t)\,\Bigl|\int_{[-2,0]}e^{i\omega\Psi(t)}\,d\Phi(\omega)\Bigr|^2
\cdot\operatorname{sgn}(\vartheta'(t)),
\end{align}
where the second line uses $\Psi' = \mu|\vartheta'|$ and
$j_{\mathrm{smooth}} = -\mu\vartheta'$ from Theorem~\ref{thm:smooth}.

\end{document}
